% GENERATED FILE. Edit canonical lessons, then run npm run generate:books.
% canonical-source-hash: fnv1a64:5db10c533ac291c2
% canonical-lessons: TE-C446-egududigudu, TE-C446-kudupu, TE-C446-tappincu, TE-C446-boltapadu, TE-C446-sottapadu

\chapter{Bumps on the Road}
\label{ch:te-bumps-on-the-road}
\hlchaptermodality{446}

\begin{chapteropening}
\textbf{By the end of this chapter:} \emph{I can say bumpy, a jolt, to avoid, to overturn and to get dented.}

Complete the last lesson of chapter 446: I can say bumpy, a jolt, to avoid, to overturn and to get dented.
\end{chapteropening}

\section[\texorpdfstring{\te{ఎగుడుదిగుడు} (eguḍudiguḍu)}{ఎగుడుదిగుడు (eguḍudiguḍu)}]{\te{ఎగుడుదిగుడు} (eguḍudiguḍu) — bumpy; uneven; ups and downs}

\label{lesson:TE-C446-egududigudu}



\begin{quote}
Before the new one: say the Telugu for to slow down, then the Telugu for to lose control.
\end{quote}

\subsection*{You'll want to know: \te{ఎగుడుదిగుడు}}
\textbf{\te{ఎగుడుదిగుడు}} — \emph{eguḍudiguḍu} — \textquotedblleft{}bumpy; uneven; ups and downs\textquotedblright{}.

\begin{quote}
\te{పల్లెకు} \te{వెళ్ళే} \te{రోడ్డు} \te{ఎగుడుదిగుడుగా} \te{ఉంది}.
\end{quote}

\begin{quote}
\emph{palleku veḷḷē rōḍḍu eguḍudiguḍugā undi.}
\end{quote}

\begin{quote}
The road to the village is bumpy.
\end{quote}

With \te{ఉండు} it takes -\te{గా}: \te{ఎగుడుదిగుడుగా} \te{ఉంది}, it is bumpy. It is also said of a life full of ups and downs.

\begin{cousinweb}
\textbf{\te{ఎగుడుదిగుడు}} (\emph{eguḍudiguḍu}) is native Dravidian, a rhyming pair: \te{ఎగుడు}, \textquotedblleft{}high, raised\textquotedblright{}, which DEDR 3730 lists beside \emph{ega diga}, \textquotedblleft{}up and down\textquotedblright{}, and \te{దిగుడు}, \textquotedblleft{}a going down\textquotedblright{}, from \te{దిగు}, \textquotedblleft{}to go down\textquotedblright{} (DEDR 502).
\end{cousinweb}

\begin{grammarlens}[title={What you've built}]
One more word for getting around.
\end{grammarlens}

\subsection*{Your turn}
\begin{itemize}
\raggedright
  \item \emph{Say it:} \emph{eguḍudiguḍu}
  \item \emph{Say it:} \emph{eguḍudiguḍu}, once more
  \item \emph{Say it:} \emph{palleku veḷḷē rōḍḍu eguḍudiguḍugā undi.} — The road to the village is bumpy.
  \item \emph{Recall:} say the Telugu for to slow down, then the Telugu for to lose control, then say \emph{eguḍudiguḍu} again
\end{itemize}

\begin{tcolorbox}[breakable,colback=teal!4,colframe=teal!35!black,title={Before you move on}]
What does \te{ఎగుడుదిగుడు} mean? (\textquotedblleft{}Bumpy; uneven; ups and downs\textquotedblright{}.) Say it once more.
\end{tcolorbox}
\section[\texorpdfstring{\te{కుదుపు} (kudupu)}{కుదుపు (kudupu)}]{\te{కుదుపు} (kudupu) — a jolt; a shake}

\label{lesson:TE-C446-kudupu}



\begin{quote}
Before the new one: say the Telugu for to lose control, then the Telugu for bumpy.
\end{quote}

\subsection*{You'll want to know: \te{కుదుపు}}
\textbf{\te{కుదుపు}} — \emph{kudupu} — \textquotedblleft{}a jolt; a shake\textquotedblright{}.

\begin{quote}
\te{బస్సు} \te{గుంటలో} \te{పడినప్పుడు} \te{పెద్ద} \te{కుదుపు} \te{వచ్చింది}.
\end{quote}

\begin{quote}
\emph{bassu guṇṭalō paḍinappuḍu pedda kudupu vaccindi.}
\end{quote}

\begin{quote}
When the bus hit the pothole there was a big jolt.
\end{quote}

The same word is a verb, to shake something: \te{సీసాను} \te{కుదుపు}, to shake a bottle.

\begin{cousinweb}
\textbf{\te{కుదుపు}} (\emph{kudupu}) is native Dravidian: DEDR 1705 gives Telugu \emph{kudupu}, \textquotedblleft{}to shake, to jolt\textquotedblright{}, and as a noun \textquotedblleft{}a jolting\textquotedblright{}, in the family of Tamil \emph{kuti}, \textquotedblleft{}to leap\textquotedblright{}.
\end{cousinweb}

\begin{grammarlens}[title={What you've built}]
One more word for getting around.
\end{grammarlens}

\subsection*{Your turn}
\begin{itemize}
\raggedright
  \item \emph{Say it:} \emph{kudupu}
  \item \emph{Say it:} \emph{kudupu}, once more
  \item \emph{Say it:} \emph{bassu guṇṭalō paḍinappuḍu pedda kudupu vaccindi.} — When the bus hit the pothole there was a big jolt.
  \item \emph{Recall:} say the Telugu for to lose control, then the Telugu for bumpy, then say \emph{kudupu} again
\end{itemize}

\begin{tcolorbox}[breakable,colback=teal!4,colframe=teal!35!black,title={Before you move on}]
What does \te{కుదుపు} mean? (\textquotedblleft{}A jolt; a shake\textquotedblright{}.) Say it once more.
\end{tcolorbox}
\section[\texorpdfstring{\te{తప్పించు} (tappiñcu)}{తప్పించు (tappiñcu)}]{\te{తప్పించు} (tappiñcu) — to avoid; to dodge}

\label{lesson:TE-C446-tappincu}



\begin{quote}
Before the new one: say the Telugu for bumpy, then the Telugu for a jolt.
\end{quote}

\subsection*{You'll want to know: \te{తప్పించు}}
\textbf{\te{తప్పించు}} — \emph{tappiñcu} — \textquotedblleft{}to avoid; to dodge\textquotedblright{}.

\begin{quote}
\te{డ్రైవరు} \te{గుంటను} \te{తప్పించి} \te{వెళ్ళాడు}.
\end{quote}

\begin{quote}
\emph{ḍraivaru guṇṭanu tappiñci veḷḷāḍu.}
\end{quote}

\begin{quote}
The driver steered round the pothole.
\end{quote}

Past (I …): \te{తప్పించాను}, \emph{tappiñcānu}.

\te{తప్పించుకో}, to escape, is the same verb with -\te{కో}: you get yourself out of something. \te{తప్పించు} is to steer clear of a thing.

\begin{cousinweb}
\textbf{\te{తప్పించు}} (\emph{tappiñcu}) is native Dravidian: DEDR 3071 lists Telugu \emph{tappincu} beside \emph{tappu}, \textquotedblleft{}to slip, to miss, to escape\textquotedblright{}.
\end{cousinweb}

\begin{grammarlens}[title={What you've built}]
One more word for getting around.
\end{grammarlens}

\subsection*{Your turn}
\begin{itemize}
\raggedright
  \item \emph{Say it:} \emph{tappiñcu}
  \item \emph{Say it:} \emph{tappiñcu}, once more
  \item \emph{Say it:} \emph{ḍraivaru guṇṭanu tappiñci veḷḷāḍu.} — The driver steered round the pothole.
  \item \emph{Recall:} say the Telugu for bumpy, then the Telugu for a jolt, then say \emph{tappiñcu} again
\end{itemize}

\begin{tcolorbox}[breakable,colback=teal!4,colframe=teal!35!black,title={Before you move on}]
What does \te{తప్పించు} mean? (\textquotedblleft{}To avoid; to dodge\textquotedblright{}.) Say it once more.
\end{tcolorbox}
\section[\texorpdfstring{\te{బోల్తాపడు} (bōltāpaḍu)}{బోల్తాపడు (bōltāpaḍu)}]{\te{బోల్తాపడు} (bōltāpaḍu) — to overturn; to topple over}

\label{lesson:TE-C446-boltapadu}



\begin{quote}
Before the new one: say the Telugu for a jolt, then the Telugu for to avoid.
\end{quote}

\subsection*{You'll want to know: \te{బోల్తాపడు}}
\textbf{\te{బోల్తాపడు}} — \emph{bōltāpaḍu} — \textquotedblleft{}to overturn; to topple over\textquotedblright{}.

\begin{quote}
\te{మలుపు} \te{దగ్గర} \te{ఒక} \te{ఆటో} \te{బోల్తాపడింది}.
\end{quote}

\begin{quote}
\emph{malupu daggara oka āṭō bōltāpaḍindi.}
\end{quote}

\begin{quote}
An auto overturned at the bend.
\end{quote}

Past (I …): \te{బోల్తాపడ్డాను}, \emph{bōltāpaḍḍānu}.

Said of a vehicle that rolls on its side. Spoken Telugu also uses it of a plan that falls through.

\begin{cousinweb}
\textbf{\te{బోల్తాపడు}} (\emph{bōltāpaḍu}) is a Telugu compound: \te{బోల్తా}, \textquotedblleft{}a tumble\textquotedblright{}, whose origin is not settled, and the native \te{పడు}, \textquotedblleft{}to fall into, to undergo\textquotedblright{} (DEDR 3853).
\end{cousinweb}

\begin{grammarlens}[title={What you've built}]
One more word for getting around.
\end{grammarlens}

\subsection*{Your turn}
\begin{itemize}
\raggedright
  \item \emph{Say it:} \emph{bōltāpaḍu}
  \item \emph{Say it:} \emph{bōltāpaḍu}, once more
  \item \emph{Say it:} \emph{malupu daggara oka āṭō bōltāpaḍindi.} — An auto overturned at the bend.
  \item \emph{Recall:} say the Telugu for a jolt, then the Telugu for to avoid, then say \emph{bōltāpaḍu} again
\end{itemize}

\begin{tcolorbox}[breakable,colback=teal!4,colframe=teal!35!black,title={Before you move on}]
What does \te{బోల్తాపడు} mean? (\textquotedblleft{}To overturn; to topple over\textquotedblright{}.) Say it once more.
\end{tcolorbox}
\section[\texorpdfstring{\te{సొట్టపడు} (soṭṭapaḍu)}{సొట్టపడు (soṭṭapaḍu)}]{\te{సొట్టపడు} (soṭṭapaḍu) — to get dented}

\label{lesson:TE-C446-sottapadu}



\begin{quote}
Before the new one: say the Telugu for to avoid, then the Telugu for to overturn.
\end{quote}

\subsection*{You'll want to know: \te{సొట్టపడు}}
\textbf{\te{సొట్టపడు}} — \emph{soṭṭapaḍu} — \textquotedblleft{}to get dented\textquotedblright{}.

\begin{quote}
\te{చిన్న} \te{ప్రమాదంలో} \te{మా} \te{కారు} \te{తలుపు} \te{సొట్టపడింది}.
\end{quote}

\begin{quote}
\emph{cinna pramādaṁlō mā kāru talupu soṭṭapaḍindi.}
\end{quote}

\begin{quote}
In a small accident the door of our car got dented.
\end{quote}

Past (it …): \te{సొట్టపడింది}, \emph{soṭṭapaḍindi}.

\te{సొట్ట}, a dent, + \te{పడు}. A dented pot is a \te{సొట్ట} \te{గిన్నె}. Its subject is a thing, so the past shown is the third person: \te{సొట్టపడింది}, it got dented.

\begin{cousinweb}
\textbf{\te{సొట్టపడు}} (\emph{soṭṭapaḍu}) is native Dravidian: \te{సొట్ట}, \textquotedblleft{}a dent, a crookedness\textquotedblright{} (DEDR 2838), with \te{పడు}, \textquotedblleft{}to undergo\textquotedblright{} (DEDR 3853).
\end{cousinweb}

\begin{grammarlens}[title={What you've built}]
One more word for getting around.
\end{grammarlens}

\subsection*{Your turn}
\begin{itemize}
\raggedright
  \item \emph{Say it:} \emph{soṭṭapaḍu}
  \item \emph{Say it:} \emph{soṭṭapaḍu}, once more
  \item \emph{Say it:} \emph{cinna pramādaṁlō mā kāru talupu soṭṭapaḍindi.} — In a small accident the door of our car got dented.
  \item \emph{Recall:} say the Telugu for to avoid, then the Telugu for to overturn, then say \emph{soṭṭapaḍu} again
\end{itemize}

\begin{tcolorbox}[breakable,colback=teal!4,colframe=teal!35!black,title={Before you move on}]
What does \te{సొట్టపడు} mean? (\textquotedblleft{}To get dented\textquotedblright{}.) Say it once more.
\end{tcolorbox}
